% 型6：誤答分析型
% よくある誤答を先に示し，どこで論理が切れているかを特定してから正答へ進む．
% 「解けたつもり」を崩す用途に向く．差がつく単元に絞って使う．
\documentclass[lualatex,paper=b5j,fontsize=10pt,twoside]{jlreq}

\usepackage{surikumi-mondaishu}

\MathMagazineSetup{
  feature={},
  series={誤答研究／数学 I},
  series-position=left,
  pattern=zigzags,
  title={絶対値と平方根},
  author={入試基礎講座},
  author-font=mincho,
  author-position=right,
  title-fill=black,
  deck={答だけを見れば正しく見える誤答を集めた．どの行で論理が切れているかを指摘できるようにする．}
}

\begin{document}

\MakeMathMagazineTitle

\begin{mmproblem}[1]{}
$\sqrt{(x-3)^2}=3-x$ を満たす実数 $x$ の範囲を求めよ．
\end{mmproblem}

\begin{msmistake}
$\sqrt{(x-3)^2}=x-3$ であるから，与式は $x-3=3-x$ となる．
これを解いて $2x=6$ すなわち $x=3$ である．
\end{msmistake}

\MSWhereWrong{$\sqrt{A^2}=A$ としている点である．正しくは $\sqrt{A^2}=|A|$
であり，$A$ の符号によって $A$ と $-A$ のどちらになるかが変わる．}

\begin{mmsolution}
$\sqrt{(x-3)^2}=|x-3|$ であるから，与式は
\[
  |x-3|=3-x
\]
となる．一般に $|A|=-A$ となるのは $A\leq 0$ のときであり，
ここでは $3-x=-(x-3)$ であるから，条件は $x-3\leq 0$ である．
よって
\[
  x\leq 3
\]
となる．
\end{mmsolution}

\MMNote{$x=3$ も解の一つだが，解は $x=3$ だけではない．誤答は「解を一つ見つけて満足する」形になっている．}

\begin{mmproblem}[2]{}
$\sqrt{x^2}+\sqrt{(x-2)^2}=2$ を満たす実数 $x$ の範囲を求めよ．
\end{mmproblem}

\begin{msmistake}
$\sqrt{x^2}=x$，$\sqrt{(x-2)^2}=x-2$ であるから，
$x+(x-2)=2$ より $x=2$ である．
\end{msmistake}

\MSWhereWrong{ここでも絶対値を外す向きを決めずに進めている．
$|x|$ と $|x-2|$ は符号の変わり目が $x=0$ と $x=2$ の二か所にあるため，
場合分けは三通りになる．}

\begin{mmsolution}
与式は $|x|+|x-2|=2$ である．

\MMCaseHeading{1}{$x<0$ のとき}
$|x|=-x$，$|x-2|=2-x$ であるから左辺は $2-2x$ となる．
$2-2x=2$ より $x=0$ であり，$x<0$ を満たさない．

\MMCaseHeading{2}{$0\leq x\leq 2$ のとき}
$|x|=x$，$|x-2|=2-x$ であるから左辺は $2$ となり，等式は常に成り立つ．

\MMCaseHeading{3}{$x>2$ のとき}
$|x|=x$，$|x-2|=x-2$ であるから左辺は $2x-2$ となる．
$2x-2=2$ より $x=2$ であり，$x>2$ を満たさない．

以上より
\[
  0\leq x\leq 2
\]
となる．
\end{mmsolution}

\MMNote{$|x|+|x-2|$ は数直線上で $0$ と $2$ への距離の和を表す．この見方をすると，答が区間になることが図から読める．}

\end{document}
